\begin{textsample}{POS Dim 3 – rn – Score 16.00 – rn/rn\_inf\_24.txt}  \label{ex:f3_pos_056}
CAMPO DE EXPERIÊNCIAS “ CORPO, \textbf{GESTOS} E MOVIMENTOS ” Direitos de aprendizagem e desenvolvimento CONVIVER com \textbf{crianças} e \textbf{adultos} experimentando marcas da cultura corporal nos \textbf{cuidados} pessoais, na dança, música, teatro, artes circenses, \textbf{escuta} de histórias e \textbf{brincadeiras}. \textbf{BRINCAR} utilizando \textbf{criativamente} o repertório da cultura corporal e do movimento.
EXPLORAR amplo repertório de movimentos, \textbf{gestos}, olhares, produção de \textbf{sons} e de \textbf{mímicas}, descobrindo modos de ocupação e de uso do espaço com o corpo.
PARTICIPAR de atividades que envolvem práticas corporais, desenvolvendo autonomia para cuidar de si.
EXPRESSAR corporalmente emoções e representações tanto nas relações cotidianas como nas \textbf{brincadeiras}, dramatizações, danças, músicas, contação de histórias.
CONHECER -SE nas diversas oportunidades de interações e explorações com seu corpo. ( BRASIL, 2018 ) DIREITOS E OBJETIVOS Pela linguagem corporal, as \textbf{crianças} se expressam, experimentam e interagem com o mundo.
Através do corpo, elas produzem conhecimentos sobre as pessoas, sobre si mesmas, e sobre o universo social e cultural, ampliando \textbf{progressivamente} a consciência dos seus limites tanto para se proteger, quanto para proteger o corpo do outro, percebendo sua completude, bem como construindo sua autonomia e segurança emotiva.
É com o corpo e suas potencialidades expressivas e comunicativas que as \textbf{crianças} também estabelecem relações, se expressam, se comunicam através da \textbf{mímica} e movimentos, \textbf{brincando}, transformando os espaços, os objetos e os elementos da natureza.
As experiências com os movimentos possibilitam a articulação entre as várias linguagens ( gestual, verbal, \textbf{plástica}, \textbf{dramática} e musical ), alternância das palavras e \textbf{gestos} ( movimentos impulsivos ou intencionais, coordenados ou espontâneos ), produção e apreciação da música e \textbf{acompanhamento} de narrações.
Também, por meio da música, do teatro, da dança e das \textbf{brincadeiras}, as \textbf{crianças} envolvem o corpo, a emoção e a linguagem, conhecendo e reconhecendo as sensações e possibilidades de seu corpo.
Na Educação \textbf{Infantil}, a centralidade do corpo ganha destaque através das atividades planejadas pelo professor, as quais possibilitam ricas e prazerosas situações em que as \textbf{crianças}, de forma \textbf{lúdica}, interagem com os pares, exploram e vivenciam experiências, ampliando o seu repertório de movimentos, \textbf{gestos}, olhares, \textbf{sons} e \textbf{mímicas} com o corpo, com vistas a identificar as várias maneiras de usá -lo.
As oportunidades dadas às \textbf{crianças}, nesse sentido, oportunizam a construção da autonomia e liberdade, de modo que consigam, com segurança, ocupar os espaços a elas apresentados.
Para isso, o trabalho pedagógico precisa considerar as especificidades de cada \textbf{criança}, respeitando suas possibilidades, interações e modos de ver a si e ao mundo, assim, é necessário fugir de ações padronizadas ou mecânicas de \textbf{gestos} e movimentos, bem como a retração de suas maneiras de se expressar corporalmente.

% matched lemmas: acompanhamento, adulto, brincadeira, brincar, criança, criativamente, cuidado, dramático, escuta, gesto, infantil, lúdico, mímico, plástico, progressivamente, som
\end{textsample}
